\chapter{Калибровочная инвариантность и движение заряженной частицы}

\lectureoverview{Магнитное поле имеет нулевую дивергенцию, поэтому его можно представить с помощью векторного потенциала. Это представление не единственно, однако именно оно позволяет включить магнитные силы в принцип действия. Мы выводим силу Лоренца, различаем канонический и механический импульсы, строим гамильтониан и используем две разные калибровки, чтобы решить задачу о движении в однородном магнитном поле. Лекция завершается рассмотрением того, как эта механика частицы вписывается в более широкий круг вопросов теории Максвелла, магнитных монополей и динамики фазового пространства.}

\section{Последний шаг в курсе механики}

Лекция начинается с небольшой неожиданности в расписании. План состоит в том, чтобы закончить тему заряженных частиц в магнитных полях, возможно добавить электрические поля, а затем решить, что останется для еще одного занятия. Слушатели замечают, что это уже десятая лекция. В таком случае задача становится более четкой: закончить магнитную задачу и с ее помощью довести до завершения формальную структуру классической механики.

Начнем с краткого повторения. Магнитное поле — векторное поле,
\begin{equation}
  \mathbf B=\mathbf B(\mathbf x),
  \qquad
  B_i(\mathbf x)=(B_x,B_y,B_z),
\end{equation}
имеющее в каждой точке модуль и направление. В этой лекции предполагается, что внешние электрическое и магнитное поля статичны. Это ограничение вводится ради ясности, а не потому, что электромагнитные поля всегда должны быть независимы от времени.

Есть и еще одно ограничение:
\begin{equation}
  \nabla\cdot\mathbf B=0.
\end{equation}
В аналогии с течением жидкости это означает, что нет концов труб, из которых магнитный поток выливается наружу, и нет стоков, в которых он исчезает. В обычной классической электродинамике нет магнитных источников и стоков. Эксперимент не установил существование элементарного магнитного монополя, хотя в лекции неоднократно оставляется открытой возможность того, что в более полной теории он может существовать.

\begin{figure}[htbp]
\centering
\includegraphics[width=0.78\textwidth]{lecture_10_figure_02.png}
\caption{Компоненты поля, условие отсутствия источников \(\nabla\cdot\mathbf B=0\) и представление через векторный потенциал появляются вместе в начале лекции.}
\lectureframe{00:04:56}
\end{figure}

Вместо того чтобы каждый раз проверять дивергенцию, когда предлагается поле, можно автоматически удовлетворить этому ограничению. Запишем
\begin{equation}
  \boxed{\mathbf B=\nabla\times\mathbf A}.
\end{equation}
Тогда тождество
\begin{equation}
  \nabla\cdot(\nabla\times\mathbf A)=0
\end{equation}
автоматически обеспечивает его выполнение. Новое поле \(\mathbf A\) — это векторный потенциал.

\section{Цена автоматического выполнения ограничения}

Введение \(\mathbf A\) решает одну проблему и создает другую: этот объект не единственен. Для любой скалярной функции \(S(\mathbf x)\),
\begin{equation}
  \mathbf A'=\mathbf A+\nabla S
\end{equation}
имеет тот же ротор, поскольку
\begin{equation}
  \nabla\times(\nabla S)=0.
\end{equation}
Следовательно, \(\mathbf A\) и \(\mathbf A'\) описывают одно и то же магнитное поле. На достаточно простой области это и есть вся локальная неоднозначность: если у двух векторных потенциалов одинаковый ротор, то их разность локально является градиентом. Глобальная топология может потребовать нескольких областей или допускать дополнительные тонкости, но для рассматриваемых здесь полей ничего этого не нужно.

Конкретный представитель называется \emph{калибровкой}, а переход от одного представителя к другому — \emph{калибровочным преобразованием}. Поле \(\mathbf B\) не меняется, поэтому оно калибровочно-инвариантно. Рабочий принцип состоит в том, что наблюдаемая физика не должна зависеть от того, какой представитель мы выбрали. Мы расширили описание, введя \(\mathbf A\), а затем устранили избыточность этого описания, потребовав калибровочно-инвариантных предсказаний.

\subsection{Три калибровки для одного однородного поля}

Лекция сразу делает эту идею конкретной. Выберем координаты так, чтобы однородное магнитное поле было направлено вдоль оси \(z\):
\begin{equation}
  \mathbf B=B\hat{\mathbf z},
  \qquad
  B_x=B_y=0,
  \qquad
  B_z=B,
\end{equation}
где \(B\) — константа. Первый векторный потенциал имеет вид
\begin{equation}
  \mathbf A^{(1)}=(0,Bx,0).
\end{equation}
Действительно,
\begin{equation}
  (\nabla\times\mathbf A^{(1)})_z
  =
  \frac{\partial A_y^{(1)}}{\partial x}
  -
  \frac{\partial A_x^{(1)}}{\partial y}
  =B.
\end{equation}
Второй выбор таков:
\begin{equation}
  \mathbf A^{(2)}=(-By,0,0),
\end{equation}
для него
\begin{equation}
  (\nabla\times\mathbf A^{(2)})_z
  =
  0-\frac{\partial(-By)}{\partial y}
  =B.
\end{equation}

Чтобы перейти между ними, возьмем
\begin{equation}
  S=Cxy,
  \qquad
  \nabla S=(Cy,Cx,0).
\end{equation}
Начиная с \(\mathbf A^{(1)}\), после преобразования получаем компоненты
\begin{equation}
  A_x'=Cy,
  \qquad
  A_y'=(B+C)x,
  \qquad
  A_z'=0.
\end{equation}
Выбор \(C=-B\) дает \(\mathbf A^{(2)}\). Выбор \(C=-B/2\) дает симметричную калибровку,
\begin{equation}
  \mathbf A^{(\mathrm{sym})}
  =
  \left(-\frac{B}{2}y,\frac{B}{2}x,0\right)
  =
  \frac12\,\mathbf B\times\mathbf r.
\end{equation}
У всех трех один и тот же ротор. Позже мы намеренно вернемся к первым двум, потому что каждая из них явно выявляет разную трансляционную симметрию.

\begin{classroomqa}
\audiencequestion{Может ли векторный потенциал обладать той же симметрией, что и однородное магнитное поле?}
\lecturerresponse{Он не может быть пространственно постоянным так, как магнитное поле, потому что у постоянного векторного потенциала ротор равен нулю. Одна калибровка также не может сделать явными все трансляционные симметрии ненулевого однородного поля. Симметричная калибровка действительно делает явными вращения вокруг оси \(z\), так что точный вывод состоит не в том, что всякая симметрия должна отсутствовать у \(\mathbf A\), а в том, что потенциал не обязан одновременно проявлять все симметрии физического поля.}
\end{classroomqa}

Семейство \(S=Cxy\) — лишь пример. В более общем случае можно наложить полезные калибровочные условия, например потребовать, чтобы в подходящей области одна компонента \(\mathbf A\) обращалась в нуль. Название \emph{калибровка} исторически возникло из неудачной попытки интерпретировать эту свободу как локальное изменение единиц измерения. От этой интерпретации отказались; название осталось.

\begin{classroomqa}
\audiencequestion{Всегда ли калибровочное преобразование означает добавление градиента скаляра?}
\lecturerresponse{Именно такое преобразование используется здесь для статического электромагнитного векторного потенциала. В других калибровочных теориях преобразования устроены сложнее, а в гравитации смену координат часто описывают как калибровочное изменение, потому что физическая геометрия не может зависеть от координатных меток. Однако в этой главе под калибровочным преобразованием понимается \(\mathbf A\mapsto\mathbf A+\nabla S\).}
\end{classroomqa}

\begin{classroomqa}
\audiencequestion{Может ли одна калибровка быть удобнее другой с вычислительной точки зрения?}
\lecturerresponse{Очень часто. Симметрия или закон сохранения могут быть явными в одной калибровке и скрытыми в другой. Это не просто абстрактная возможность: позже в этой лекции одна калибровка сделает \(y\) циклической координатой, а другая — \(x\). Оба вычисления описывают одну и ту же орбиту.}
\end{classroomqa}

\begin{classroomqa}
\audiencequestion{Связана ли эта неоднозначность с выбором собственных подпространств или нахождением собственных значений?}
\lecturerresponse{Прямой связи, которая была бы здесь нужна, нет. В обеих темах приходится делать полезный математический выбор, но переход к другому собственному базису и электромагнитное калибровочное преобразование — это разные структуры.}
\end{classroomqa}

\section{Магнитный лагранжиан}

Теперь вернемся к механике. Объект обсуждения — лагранжиан точечной частицы, движущейся в заданном магнитном поле, а не еще лагранжиан самого электромагнитного поля. В используемой на лекции системе единиц
\begin{equation}
  L
  =
  \frac{m}{2}
  \left(\dot x^2+\dot y^2+\dot z^2\right)
  +
  \frac{e}{c}
  \left(
    A_x\dot x+A_y\dot y+A_z\dot z
  \right).
\end{equation}
В векторных обозначениях,
\begin{equation}
  \boxed{
  L=\frac12 m\mathbf v^2+\frac{e}{c}\mathbf A(\mathbf x)\cdot\mathbf v
  }.
\end{equation}
Коэффициент \(e\) — это заряд частицы. Явное присутствие \(c\) зависит от выбора единиц; структурно существенным является зависящее от скорости взаимодействие \(\mathbf A\cdot\mathbf v\).

\begin{figure}[htbp]
\centering
\includegraphics[width=0.86\textwidth]{lecture_10_figure_03.png}
\caption{Лагранжиан частицы в магнитном поле записан покомпонентно, так что все три скоростных члена связи остаются видимыми.}
\lectureframe{00:25:29}
\end{figure}

Сначала этот член постулируется, а не выводится. Нам предлагают попробовать его и посмотреть, какие уравнения он порождает. Можно было бы также попробовать член, построенный непосредственно из \(\mathbf B\), но он дал бы производные от \(\mathbf B\), а не силу Лоренца. Векторный потенциал — локальный объект, который правильно связывается с траекторией.

\subsection{Почему калибровочный сдвиг не меняет траекторию}

Прежде чем варьировать действие, проверим, что делает с ним калибровочное преобразование. При фиксированных начальной и конечной пространственно-временных точках,
\begin{equation}
  \mathcal S[\mathbf x]
  =
  \int_{t_i}^{t_f}L\,dt.
\end{equation}
При \(\mathbf A\mapsto\mathbf A+\nabla S\) изменение равно
\begin{align}
  \Delta\mathcal S
  &=
  \frac{e}{c}
  \int_{t_i}^{t_f}
  \nabla S(\mathbf x)\cdot\dot{\mathbf x}\,dt\\
  &=
  \frac{e}{c}
  \int_{t_i}^{t_f}\frac{dS(\mathbf x(t))}{dt}\,dt\\
  &=
  \frac{e}{c}
  \left[
    S(\mathbf x_f)-S(\mathbf x_i)
  \right].
\end{align}
Следовательно, численное значение действия меняется. Но добавленный член зависит только от фиксированных концов, а не от варьируемой внутренней части траектории. Его вариация равна нулю, и он не может изменить стационарную траекторию.

\begin{classroomqa}
\audiencequestion{Что значит утверждение, что уравнения калибровочно-инвариантны?}
\lecturerresponse{После применения уравнений Эйлера--Лагранжа \(\mathbf A\) должен появляться только в роторных комбинациях, распознаваемых как компоненты \(\mathbf B\). Сам потенциал может меняться, и действие может сдвигаться на концевой член, но физическая траектория и закон силы должны оставаться неизменными.}
\end{classroomqa}

\section{Два понятия импульса}

Уравнения Эйлера--Лагранжа можно записать в виде
\begin{equation}
  \dot p_i=\frac{\partial L}{\partial x_i},
  \qquad
  p_i=\frac{\partial L}{\partial\dot x_i}.
\end{equation}
Здесь нужно различать два импульса. Привычный механический импульс равен
\begin{equation}
  \boldsymbol{\pi}=m\mathbf v.
\end{equation}
Канонический импульс определяется лагранжианом:
\begin{equation}
  p_i
  =
  m\dot x_i+\frac{e}{c}A_i.
\end{equation}
Следовательно,
\begin{equation}
  \boxed{
  \boldsymbol{\pi}
  =
  \mathbf p-\frac{e}{c}\mathbf A
  =
  m\mathbf v
  }.
\end{equation}
Для свободной частицы, или при обычном потенциале, зависящем от положения, канонический и механический импульсы совпадают. Магнитное взаимодействие зависит от скорости, поэтому здесь они уже не совпадают.

\begin{figure}[htbp]
\centering
\includegraphics[width=0.78\textwidth]{lecture_10_figure_04.png}
\caption{На доске каноническое определение \(p_i=\partial L/\partial\dot x_i\) отделено от \(m\dot x_i\), после чего начинают дифференцировать это выражение вдоль траектории.}
\lectureframe{00:37:33}
\end{figure}

Это различие — не просто вопрос терминологии. Канонический импульс — это величина, входящая в уравнения Гамильтона, скобки Пуассона и закон сохранения, связанный с циклической координатой. Механический импульс — это непосредственно измеряемое произведение массы на скорость.

\begin{classroomqa}
\audiencequestion{Является ли \((e/c)\mathbf A\cdot\mathbf v\) кинетической или потенциальной энергией, и как понять, нужно ли его прибавлять или вычитать?}
\lecturerresponse{Это не обычная кинетическая энергия и не обычная зависящая от положения потенциальная энергия. Это зависящее от скорости взаимодействие. На данном этапе его знак является частью предлагаемой связи; изменение этого знака эквивалентно изменению знака заряда или \(\mathbf A\). Решающая проверка — это следующий отсюда закон силы.}
\end{classroomqa}

\section{Выведение силы Лоренца из алгебры}

В лекции одна компонента разбирается медленно, шаг за шагом. Для \(p_x\),
\begin{equation}
  p_x=m\dot x+\frac{e}{c}A_x(\mathbf x).
\end{equation}
Хотя внешнее поле статично, \(A_x(\mathbf x(t))\) меняется по мере движения частицы. Его полная производная по времени равна
\begin{equation}
  \frac{dA_x}{dt}
  =
  \frac{\partial A_x}{\partial x}\dot x
  +
  \frac{\partial A_x}{\partial y}\dot y
  +
  \frac{\partial A_x}{\partial z}\dot z.
\end{equation}
Следовательно,
\begin{equation}
  \dot p_x
  =
  m\ddot x
  +
  \frac{e}{c}
  \left(
    \partial_xA_x\,\dot x
    +
    \partial_yA_x\,\dot y
    +
    \partial_zA_x\,\dot z
  \right).
\end{equation}
Правая часть уравнения Эйлера--Лагранжа равна
\begin{equation}
  \frac{\partial L}{\partial x}
  =
  \frac{e}{c}
  \left(
    \partial_xA_x\,\dot x
    +
    \partial_xA_y\,\dot y
    +
    \partial_xA_z\,\dot z
  \right).
\end{equation}
Члены \(\partial_xA_x\,\dot x\) сокращаются. Собирая оставшееся по компонентам скорости, получаем
\begin{align}
  m\ddot x
  &=
  \frac{e}{c}
  \left[
    \left(\partial_xA_y-\partial_yA_x\right)\dot y
    +
    \left(\partial_xA_z-\partial_zA_x\right)\dot z
  \right].
\end{align}

\begin{figure}[htbp]
\centering
\includegraphics[width=0.80\textwidth]{lecture_10_figure_05.png}
\caption{Завершённое вычисление \(x\)-компоненты: производные векторного потенциала объединились в роторные комбинации.}
\lectureframe{00:46:41}
\end{figure}

Используя
\begin{align}
  B_y&=\partial_zA_x-\partial_xA_z,\\
  B_z&=\partial_xA_y-\partial_yA_x,
\end{align}
получаем
\begin{equation}
  m a_x
  =
  \frac{e}{c}
  \left(B_zv_y-B_yv_z\right)
  =
  \frac{e}{c}
  (\mathbf v\times\mathbf B)_x.
\end{equation}
Тот же расчёт для \(y\) и \(z\) даёт
\begin{equation}
  \boxed{
  m\mathbf a=\frac{e}{c}\,\mathbf v\times\mathbf B
  }.
\end{equation}
Это закон магнитной силы Лоренца. В лекции полное уравнение называется уравнением Ньютона--Лоренца: Ньютон даёт \(m\mathbf a\), Лоренц — магнитную силу.

Результат проходит все три проверки, поставленные до вычисления. Новый член даёт ненулевую силу. Сила калибровочно инвариантна, потому что остаётся только \(\nabla\times\mathbf A=\mathbf B\). И векторный потенциал, хотя и избыточен, дал магнитной силе принцип действия. Замена \(\mathbf A\) на \(\mathbf B\) во взаимодействии привела бы к комбинациям с ротором \(\mathbf B\) и не воспроизвела бы этот закон.

\section{Гамильтониан и смысл энергии}

Следующий вопрос состоит в том, допускает ли та же система гамильтоново описание. Начнём с
\begin{align}
  p_x&=m\dot x+\frac{e}{c}A_x,\\
  p_y&=m\dot y+\frac{e}{c}A_y,\\
  p_z&=m\dot z+\frac{e}{c}A_z.
\end{align}
Эти соотношения можно обратить:
\begin{align}
  \dot x&=\frac{1}{m}\left(p_x-\frac{e}{c}A_x\right),\\
  \dot y&=\frac{1}{m}\left(p_y-\frac{e}{c}A_y\right),\\
  \dot z&=\frac{1}{m}\left(p_z-\frac{e}{c}A_z\right).
\end{align}
Преобразование Лежандра имеет вид
\begin{equation}
  H=\sum_i p_i\dot x_i-L.
\end{equation}
Подстановка и сокращение дают
\begin{equation}
  \boxed{
  H(\mathbf x,\mathbf p)
  =
  \frac{1}{2m}
  \left(
    \mathbf p-\frac{e}{c}\mathbf A(\mathbf x)
  \right)^2
  }.
\end{equation}

\begin{figure}[htbp]
\centering
\includegraphics[width=0.86\textwidth]{lecture_10_figure_06.png}
\caption{Завершённое построение гамильтониана: слева остаются три канонических импульса, справа — квадраты минимально связанной комбинации.}
\lectureframe{00:53:05}
\end{figure}

Записанный в канонических переменных, он отличается от гамильтониана свободной частицы. Однако, записанный через физическую скорость,
\begin{equation}
  H=\frac{\boldsymbol{\pi}^{\,2}}{2m}
  =\frac12m\mathbf v^2.
\end{equation}
Противоречия нет. Одно и то же численное значение энергии представлено другой функцией от \((\mathbf x,\mathbf p)\), поэтому уравнения Гамильтона различны, хотя энергия по-прежнему равна обычной кинетической энергии.

Поскольку поле статично, \(H\) сохраняется. Следовательно, модуль скорости постоянен. Тот же факт наглядно выражает сила Лоренца:
\begin{equation}
  \frac{d}{dt}\left(\frac12mv^2\right)
  =
  m\mathbf v\cdot\mathbf a
  =
  \frac{e}{c}\,
  \mathbf v\cdot(\mathbf v\times\mathbf B)
  =0.
\end{equation}
Магнитное поле изменяет направление движения, а не модуль скорости. Оно искривляет траекторию, но работы не совершает.

\section{Однородное поле: решение с помощью смены калибровки}

После перерыва снова рассмотрим однородное поле
\begin{equation}
  \mathbf B=B\hat{\mathbf z}.
\end{equation}
Цель состоит не просто в том, чтобы решить задачу о знакомой орбите. Нужно показать, почему выбор калибровки может сделать расчёт, основанный на симметрии, почти не требующим усилий.

\subsection{Первая калибровка: координаты y и z — циклические}

Возьмём
\begin{equation}
  A_x=A_z=0,
  \qquad
  A_y=Bx.
\end{equation}
Тогда
\begin{equation}
  L
  =
  \frac{m}{2}
  \left(\dot x^2+\dot y^2+\dot z^2\right)
  +
  \frac{eB}{c}x\dot y.
\end{equation}
Координаты \(y\) и \(z\) явно не входят. Они циклические, поэтому их канонические импульсы сохраняются. Для \(z\),
\begin{equation}
  p_z=m\dot z,
  \qquad
  \dot p_z=0,
\end{equation}
и, следовательно,
\begin{equation}
  a_z=0.
\end{equation}
Частица движется вдоль поля равномерно.

Для \(y\),
\begin{equation}
  p_y=m\dot y+\frac{eB}{c}x,
  \qquad
  \dot p_y=0.
\end{equation}
Поэтому
\begin{equation}
  m a_y+\frac{eB}{c}v_x=0,
  \qquad
  \boxed{m a_y=-\frac{eB}{c}v_x}.
\end{equation}

\begin{figure}[htbp]
\centering
\includegraphics[width=0.84\textwidth]{lecture_10_figure_07.png}
\caption{В калибровке \(A_y=Bx\) на доске записаны сохраняющиеся \(p_z\) и \(p_y\) и вытекающая отсюда \(y\)-компонента силы Лоренца.}
\lectureframe{01:03:50}
\end{figure}

\subsection{Вторая калибровка: координата x — циклическая}

Теперь перейдём к
\begin{equation}
  A_x=-By,
  \qquad
  A_y=A_z=0.
\end{equation}
Лагранжиан принимает вид
\begin{equation}
  L
  =
  \frac{m}{2}
  \left(\dot x^2+\dot y^2+\dot z^2\right)
  -
  \frac{eB}{c}y\dot x.
\end{equation}
На этот раз циклической является \(x\):
\begin{equation}
  p_x=m\dot x-\frac{eB}{c}y,
  \qquad
  \dot p_x=0.
\end{equation}
Следовательно,
\begin{equation}
  \boxed{m a_x=\frac{eB}{c}v_y}.
\end{equation}
Вместе с результатом из первой калибровки и \(a_z=0\) это в точности даёт
\begin{equation}
  m\mathbf a=\frac{e}{c}\mathbf v\times B\hat{\mathbf z}.
\end{equation}

\begin{figure}[htbp]
\centering
\includegraphics[width=0.84\textwidth]{lecture_10_figure_08.png}
\caption{После смены калибровки видны обе поперечные компоненты силы Лоренца. Каждая возникает как закон сохранения канонического импульса в своей калибровке.}
\lectureframe{01:08:05}
\end{figure}

Физическая трансляционная симметрия присутствовала всё время. Выбор калибровки изменил лишь то, какая часть этой симметрии проявляется как обычная циклическая координата. Поэтому канонический импульс может сохраняться в одной калибровке и иметь другой вид в другой, тогда как измеряемое ускорение остаётся тем же.

\begin{classroomqa}
\audiencequestion{Являются ли компоненты канонического импульса калибровочно инвариантными? Являются ли они непосредственно физическими величинами?}
\lecturerresponse{Нет. При \(\mathbf A\mapsto\mathbf A+\nabla S\) канонический импульс изменяется на \((e/c)\nabla S\). Механический импульс \(m\mathbf v=\mathbf p-(e/c)\mathbf A\) калибровочно инвариантен и напрямую связан со скоростью частицы. Тем не менее канонический импульс незаменим, потому что именно эта переменная используется в уравнениях Гамильтона и скобках Пуассона.}
\end{classroomqa}

\section{Орбита в однородном магнитном поле}

Формальная часть закончена; теперь посмотрим на движение. Сначала предположим \(v_z=0\), так что частица остаётся в плоскости \(xy\). Ускорение перпендикулярно скорости, а его модуль постоянен:
\begin{equation}
  |\mathbf a|
  =
  \frac{|e|}{mc}vB.
\end{equation}
Частица, движущаяся с постоянной скоростью и испытывающая ускорение постоянного модуля, всегда перпендикулярное её скорости, движется по окружности.

Если \(r\) — радиус, то центростремительное ускорение равно \(v^2/r\). Приравнивая модули, получаем
\begin{equation}
  m\frac{v^2}{r}
  =
  \frac{|e|}{c}vB,
\end{equation}
и, следовательно,
\begin{equation}
  \boxed{
  r=\frac{mcv}{|e|B}
  }.
\end{equation}
В лекции сначала это выражение записывается без \(c\), затем замечается несоответствие с выбранной конвенцией для силы Лоренца, и этот множитель восстанавливается. Знак \(eB\) определяет направление вращения; сам радиус положителен.

\begin{figure}[htbp]
\centering
\includegraphics[width=0.84\textwidth]{lecture_10_figure_09.png}
\caption{Круговая орбита, касательная скорость и соотношение для радиуса после восстановления пропущенного множителя \(c\).}
\lectureframe{01:16:50}
\end{figure}

Поле фиксирует связь между скоростью и радиусом, но не центр окружности. Центр может находиться где угодно в поперечной плоскости. Большая поперечная скорость даёт большую орбиту; больший заряд или более сильное магнитное поле дают орбиту меньшего радиуса.

Если \(v_z\neq0\), поперечное движение остаётся круговым, тогда как координата \(z\) меняется равномерно. Объединение этих двух движений даёт винтовую линию вокруг направления магнитного поля.\editorialnote{В лекции эти две составляющие выводятся по отдельности: \(a_z=0\) и круговое поперечное движение. В сочетании они и дают указанную винтовую траекторию.}

\section{Откуда взялось это взаимодействие?}

Расчёт орбиты побуждает аудиторию вернуться к началу. Был ли член \((e/c)\mathbf A\cdot\mathbf v\) просто заранее известен? Честный ответ состоит в том, что имеется несколько логически согласованных отправных точек.

\begin{classroomqa}
\audiencequestion{Неужели мы просто знали магнитный лагранжиан и записали его?}
\lecturerresponse{Для данного изложения — да: записываем линейное по скорости взаимодействие и исследуем его. Исходя из этого члена, обнаруживают калибровочную инвариантность, а затем видят, что роторные комбинации в уравнениях ведут себя как магнитное поле. Можно вместо этого начать с наблюдаемой силы Лоренца и искать действие или исходить из современного требования калибровочной инвариантности. Исторически эти элементы были открыты в менее прямом порядке.}
\end{classroomqa}

Есть и полезная проверка того, допускает ли предложенный лагранжиан обычное гамильтоново построение. Уравнения
\begin{equation}
  p_i=\frac{\partial L}{\partial\dot x_i}
\end{equation}
должны быть разрешимы относительно скоростей, по крайней мере в рассматриваемой здесь регулярной системе. Говоря точнее, гессиан по скоростям
\begin{equation}
  \frac{\partial^2L}
       {\partial\dot x_i\,\partial\dot x_j}
\end{equation}
должен быть невырожденным. Для этой частицы он просто равен \(m\delta_{ij}\), так что преобразование Лежандра однозначно.

\begin{classroomqa}
\audiencequestion{Что делает предложенный лагранжиан допустимым?}
\lecturerresponse{Для обычной механики, развиваемой здесь, он должен давать однозначные уравнения движения и допускать единственный гамильтониан после выражения скоростей через импульсы. Сингулярные лагранжианы действительно встречаются в калибровочных теориях поля, но для рассмотрения их связей нужен аппарат, выходящий за рамки этого курса.}
\end{classroomqa}

\begin{classroomqa}
\audiencequestion{Даёт ли сама калибровочная симметрия сохраняющуюся величину?}
\lecturerresponse{Не в одной лишь задаче о частице в заданном внешнем поле. Чтобы корректно обсудить соответствующий закон сохранения, электромагнитное поле должно быть включено в число динамических переменных. В полной теории поля калибровочная инвариантность связана с сохранением заряда, но для этого нужны действие поля и уравнения Максвелла.}
\end{classroomqa}

\section{Фиксированные фоны и динамика Максвелла}

До сих пор ничто не определяет, как должно эволюционировать электромагнитное поле. Поле задавалось как фиксированный фон, и вопрос состоял лишь в следующем: как движется частица при заданных \(\mathbf E\) и \(\mathbf B\)?

\begin{classroomqa}
\audiencequestion{Как это связано с уравнениями Максвелла?}
\lecturerresponse{Чтобы вывести уравнения Максвелла, нужно добавить собственное действие электромагнитного поля и варьировать не только траекторию частицы, но и само поле. Связь \((e/c)\int\mathbf A\cdot d\mathbf x\) составляет одну часть полного действия. Полевой член не влияет на вариацию траектории частицы, когда фон считается фиксированным, но он существен, когда динамическим становится само поле.}
\end{classroomqa}

Классически можно было бы пренебречь формализмом действия и просто постулировать
\begin{equation}
  m\mathbf a=\frac{e}{c}\mathbf v\times\mathbf B.
\end{equation}
Для механики пробной частицы это вполне осмысленное дифференциальное уравнение. Но лагранжевы и гамильтоновы структуры становятся незаменимыми в квантовой механике, где векторный потенциал входит непосредственно.

\section{Монополи и картина с тонким соленоидом}

Что, если
\begin{equation}
  \nabla\cdot\mathbf B\neq0?
\end{equation}
Тогда \(\mathbf B\) нельзя представить всюду одним гладким векторным потенциалом, потому что дивергенция любого ротора равна нулю. Простая глобальная конструкция, использованная в этой главе, перестает работать.

На лекции эта проблема иллюстрируется воображаемым чрезвычайно длинным, исключительно тонким стержневым магнитом или соленоидом. Вблизи одного конца, если другой унести сколь угодно далеко, внешнее поле напоминает поле одного магнитного полюса. Возвращающийся поток не исчез; он заключен внутри тонкой трубки. Если бы трубку удалось сделать идеальной и ненаблюдаемой, она играла бы роль струны, прикрепленной к кажущемуся монополю.

\begin{figure}[htbp]
\centering
\includegraphics[width=0.84\textwidth]{lecture_10_figure_10.png}
\caption{Схема на доске с тонким соленоидом: силовые линии расходятся от кажущегося полюса, тогда как возвращающийся поток заключен в узкой трубке.}
\lectureframe{01:27:35}
\end{figure}

В рамках обычного классического электромагнетизма сжатие фиксированного потока в исчезающее поперечное сечение увеличивает напряженность поля и энергию поля. На лекции намеренно используются утрированные ответы о бесконечной энергии, чтобы подчеркнуть, что реалистический монополь потребовал бы новой физики, а не просто обычного соленоида, который делают все тоньше и тоньше. В квантовой теории должным образом квантованная трубка потока может стать ненаблюдаемой вне своей сердцевины; именно на этой идее основана конструкция струны Дирака, хотя полная теория монополей выходит за рамки данного курса.

\begin{classroomqa}
\audiencequestion{У точечной частицы может быть магнитный момент. Является ли это монополем?}
\lecturerresponse{Нет. Магнитный момент — это диполь: у него есть равные и противоположные полюсные характеристики. Монополь нес бы ненулевой суммарный магнитный заряд. Если разрезать обычный магнит, его полюса не изолируются; у каждого куска появляются собственные северный и южный концы.}
\end{classroomqa}

\begin{classroomqa}
\audiencequestion{Аналогична ли трудность разделения двух концов конфайнменту кварков?}
\lecturerresponse{Это полезная аналогия. При раздвигании конфайнированных кварков в обычных наблюдениях не получается изолированного кварка; соединяющее поле порождает новые связанные объекты. Здесь не утверждается тождество механизмов, но картина конечных точек, соединенных неизбежной структурой потока, похожа.}
\end{classroomqa}

Заключение лекции осторожно. Магнитные монополи не входят в обычную классическую теорию, использованную в выводе, но они содержатся во многих теориях высоких энергий. Если они существуют как элементарные частицы, то, как ожидается, они связаны с масштабами энергии, далеко выходящими за рамки механики, изучаемой в этой аудитории.

\section{Другая алгебра Гамильтона и незавершенная тема}

Последний вопрос из аудитории касается кватернионов Гамильтона. Исторически это отдельная тема по отношению к гамильтоновой механике. Гамильтон искал алгебру, которая могла бы играть для трехмерных вращений роль, подобную той, какую комплексные числа играют для вращений на плоскости. Результат использует три мнимые единицы \(i,j,k\), причем
\begin{equation}
  i^2=j^2=k^2=ijk=-1,
\end{equation}
и, например,
\begin{equation}
  ij=k,
  \qquad
  ji=-k.
\end{equation}
Единичные кватернионы дают эффективное представление трехмерных вращений. В квантовой механике ту же алгебраическую структуру можно представить матрицами Паули с подходящими множителями обычной мнимой единицы. Эта связь реальна, но она не является источником канонического импульса или гамильтониана, развитых в этом курсе.

Последняя из запрошенных тем — хаос. Времени на ее изложение не осталось. Лекция завершается, отсылая к уже построенному аппарату: динамика — это движение вдоль орбит в фазовом пространстве, и именно этот язык служит отправной точкой для изучения регулярного и хаотического поведения. Формальная структура классической механики достаточно полна, чтобы слушатели могли продолжить отсюда.
